\documentclass[11pt]{article}
\usepackage[a4paper,margin=2.4cm]{geometry}
\usepackage{xeCJK}
\usepackage{amsmath,amssymb}
\usepackage{enumitem}
\usepackage{xcolor}
\usepackage{hyperref}
\setCJKmainfont{Microsoft YaHei}
\setlist{nosep,leftmargin=1.6em}
\pagestyle{plain}

% ── 盒子：手写，只用核心命令 ──────────────────────────────────
% 原本用的是 tcolorbox，但便携 MiKTeX 里缺它的依赖 trimspaces.sty，
% 而自动装包在这台机器上没生效。与其去补依赖，不如不用它——
% 手写盒子的好处是**只依赖 LaTeX 内核**，换任何一台装了 xelatex 的机器都能编。
%
% ⚠️ 这里刻意用 \newcommand 而不是 \newenvironment：
% \newenvironment{...}[2]{begin}{end} 的**结束段不能引用参数**，
% 写了 #1 会报 "Illegal parameter number in definition of \end...".
% 普通命令没有这个限制。
\newcommand{\hwboxout}[3]{%
  \par\medskip\noindent
  \setlength{\fboxsep}{7pt}%
  \fcolorbox{#1}{#2}{\begin{minipage}{0.90\textwidth}#3\end{minipage}}%
  \par\medskip}

\title{\textbf{作业求解结果}}
\author{由 homework-solver 自动生成（求解 $\to$ 自校验 $\to$ 排版）}
\date{\today}
\begin{document}
\maketitle
\vspace{-1.2em}
\hwboxout{gray!70}{gray!8}{\textbf{本文件由流水线自动生成。}每题都经过\textbf{自校验}——
解出的结果会代回原题复算一次；\textbf{自校验未通过的题会被红框标出}，
不会安静地混在正确答案里。}
\vspace{0.2em}

\section*{总览}

\begin{itemize}

  \item 题目总数：\textbf{2}
  \item 成功求解：\textbf{2}
  \item 自校验通过：\textbf{0}
  \item 自校验\textbf{未通过}：\textbf{0}
\end{itemize}

\section*{第 1 题\quad 应用题（一元二次建模）}

\textbf{题目.}\;一个长方形的长比宽多 3 厘米，它的面积是 40 平方厘米。求这个长方形的长和宽。


\textbf{求解过程.}

\begin{enumerate}[label=\arabic*.]

  \item $设长方形的宽为x\text{cm}$
  \item $则其长为x+3\text{cm}$
  \item $根据面积公式，有：(x+3)x=40$
  \item $即：x^2 + 3x - 40 = 0$
\end{enumerate}

\hwboxout{gray!70}{gray!8}{\textbf{答案：}$x_1 = 5, x_2 = -8$}

\textbf{自校验：}不适用——由模型求解，无自动校验；仅供人工复核

\vspace{0.6em}

\section*{第 2 题\quad 古典概型（文字题）}

\textbf{题目.}\;从 1, 2, 3, 4, 5 这五个数字中任取两个不同的数，求这两个数的和为偶数的概率。


\textbf{求解过程.}

\begin{enumerate}[label=\arabic*.]

  \item $A$ 表示第一个数是奇数，$B$ 表示第二个数也是奇数。
  \item $C$ 表示第一个数是偶数，$D$ 表示第二个数也是偶数。
  \item $P(A) = P(B) = \frac{3}{5}$，因为有 3 个奇数（1, 3, 5）和 2 个偶数（2, 4）。
  \item $P(C) = P(D) = \frac{2}{5}$，同理。
  \item 两个数的和为偶数的情况有两种：$A$ 和 $B$；$C$ 和 $D$。
  \item 因此，总概率是 $P(A 	ext{ and } B) + P(C 	ext{ and } D)$。
  \item $P(A 	ext{ and } B) = P(A) \times P(B) = \frac{3}{5} \times \frac{2}{5}$；
  \item $P(C 	ext{ and } D) = P(C) \times P(D) = \frac{2}{5} \times \frac{1}{5}$。
  \item 所以，总概率是 $\frac{3}{5} \times \frac{2}{5} + \frac{2}{5} \times \frac{1}{5}$。
  \item 总概率计算如下：$P = \frac{6}{25} + \frac{2}{25} = \frac{8}{25}$。
\end{enumerate}

\hwboxout{gray!70}{gray!8}{\textbf{答案：}\(\frac{8}{25}\)}

\textbf{自校验：}不适用——由模型求解，无自动校验；仅供人工复核

\vspace{0.6em}


\end{document}
